\begin{ledgroupsized}[r]{120mm}
\footnotesize
\pstart
\noindent\textbf{Überlieferung:}
\pend
\end{ledgroupsized}
\begin{ledgroupsized}[r]{114mm}
\footnotesize
\pstart
\parindent -6mm
\makebox[6mm][l]{\textit{L}}%
Konzept:
LH~XXXVII~5 Bl.~64\textendash65. 
Ein Bogen~2\textsuperscript{o};
\href{\Scan{LH_37_05_065wz}}%
{Wasserzeichen} in Bl.~65,
\href{\Scan{LH_37_05_064wz}}%
{Gegenmarke} in Bl.~64;
Papiererhaltungsmaßnahmen;
Ränder ausgefranst; 
geringfügiger Textverlust durch Papierschaden im oberen Bereich von Bl.~64~r\textsuperscript{o} und 65~r\textsuperscript{o}.
Dreieinhalb Seiten.
\pend
\end{ledgroupsized}
\vspace{5mm}
\begin{ledgroup}
\footnotesize
\pstart
\noindent
\textbf{Datierungsgründe:}
Das Wasserzeichen im Träger von N.\;\ref{60309} ist für den angegebenen Zeitraum belegt:
erstens in einem mit Sicherheit im September 1685 verfassten Konzept für den Aufsatz \glqq Demonstratio Geometrica regulae apud Staticos receptae\grqq\ (N.\;\ref{32313} in diesem Band);
zweitens in einem Entwurf von Januar 1687 zu Leibnizens erster Streitschrift gegen
\indexN{\textls{Catelan}, François, Abbé de gest. nach 1719}Catelan
(\cite{02086}"Replique de M.\;L. à M.\;l'abbé D.\;C.", \textit{NRL}, Februar 1687, S.\;131\textendash145).
Mangels weiterer Anhaltspunkte wird für die Entstehung von N.\;\ref{60309} diese Zeitspanne vorgeschlagen.
\pend
\pstart
Das vorliegende Konzept stellt eine ausführliche Weiterentwicklung einiger Gedanken und Beweise aus den ab 1679 entstandenen Texten zum Schwerpunkt N.\;\ref{Cc2-1421}, insbesondere der zweiten Aufzeichnung (\refpassage{35_08_30_133-teil2a}{35_08_30_133-teil2b}).
Viele Diagramme und Inhalte des genannten Stücks werden in N.\;\ref{60309} wieder aufgegriffen und es sind viele wörtliche Anklänge feststellbar, wie bspw. zwischen \refpassage{37_05_064-065_4a}{37_05_064-065_4b} von N.\;\ref{60309} und \refnote{FnC:35_08_30_133-Lamy} von N.\;\ref{Cc2-1421}.
Insbesondere die Passage \refpassage{35_08_30_133-teil2a}{35_08_30_133_4} aus der zweiten Aufzeichnung wird hier in einer Randanmerkung zu \refnote{FnA:37_05_064-065_elementa-lineae} wortgleich übernommen.
(Das Abhängigkeitsverhältnis des vorliegenden Konzept von den Aufzeichnungen N.\;\ref{Cc2-1421} wird mit als Grund für deren wahrscheinlichen Datierung bis 1687, also vor dem Endpunkt der vermutlichen Entstehungsspanne von N.\;\ref{60309}, angeführt.)
\pend 
\end{ledgroup}
\selectlanguage{latin}
\vspace{8mm}
\pstart
\centering
\edtext{%
\href{\Scan{LH_37_05_064r}}%
{\lbrack64~r\textsuperscript{o}\rbrack}%
}{\lemma{}\Afootnote{\textit{Am oberen Blattrand}: Quae prima pagina habentur certa sunt et usitata, sequentia \mbox{\textlangle\textendash\ \textendash\ \textendash\ \textendash\textrangle}}}
\hfill
Singulares quaedam ratiocinationes circa 
\edlabel{37_05_064-065_1a}%
\edtext{}{\xxref{37_05_064-065_1a}{37_05_064-065_1b}\lemma{}\Bfootnote{%
centrum \bars{gravitates \textit{ändert Hrsg.}}
\lbrack/\rbrack\
\textit{(1)}~Saepe mihi %
\textit{(2)}~Non~\textit{L}}}%
centrum gravitat\lbrack i\rbrack s
\hfill
\phantom{\lbrack64~r\textsuperscript{o}\rbrack}
\pend
\vspace{0.5\baselineskip}
\nobreak
\pstart
\noindent
\edtext{%
Non\edlabel{37_05_064-065_1b}
paucae jam olim exoticae ratiocinationes\indexS{ratiocinationes exoticae circa centrum gravitatis} circa centrum gravitatis\indexS{centrum gravitatis} mihi incidere,%
}{%
\lemma{Non \bracks{...} incidere}%
\Cfootnote{%
Womöglich sind hiermit die Aufzeichnungen N.\;\ref{Cc2-1421} gemeint.}}
his nuper accessere quae sequuntur: 
\pend
\begin{edfigure}%
\includegraphics[scale=1]{%
gesamttex/edit_VIII-4/images/LH_37_05_064-065_d1_064r.pdf}
\edcaption{\lbrack\textit{Fig.~1}\rbrack}
\end{edfigure}
\pstart 
Sit semicirculus \textit{ABD} cujus centrum \textit{C}, patet eum resolvi\indexS{resolutio semicirculi in innumera triangula} posse in innumera triangula ut $CB\beta$. Trianguli autem cujusque centrum gravitatis\indexS{centrum gravitatis trianguli} \textit{e}, una tertia altitudinis abesse a basi $B\beta$. Ergo locum omnium centrorum gravitatis\indexS{locum centrorum gravitatis triangulorum} talium triangulorum 
\edtext{fore semicircumferentiam}{\lemma{fore}\Bfootnote{%
\textit{(1)}~circumferentiam %
\textit{(2)}~semicircumferentiam~\textit{L}}}
centro \textit{C} radio \textit{Ce} vel \textit{CF}, qui est duabus tertiis radii \textit{CA} aequalis, descriptam. Centrum autem gravitatis\indexS{centrum gravitatis} omnium horum punctorum \textit{e} ex quibus pendere pondera suis triangulis 
\edtext{adeoque hoc loco}{\lemma{adeoque}\Bfootnote{%
\textit{(1)}~inter se aeq %
\textit{(2)}~hoc loco~\textit{L}}} 
in circulo inter se aequalia, intelligantur, erit centrum gravitatis totius semicirculi.\indexS{centrum gravitatis semicirculi} 
Videamus an centrum gravitatis\indexS{centrum gravitatis} omnium punctorum \textit{e}, coincidat cum centro gravitatis ipsius loci\indexS{centrum gravitatis loci punctorum seu curvae} omnium horum punctorum seu curvae. Et quidem cum puncta ista \textit{e}, sint etiam centra gravitatis\indexS{centrum gravitatis rectae} rectarum $L\lambda$, elementorum circumferentiae \textit{FeD} hinc si semper 
\edtext{contingat rectas}{\lemma{contingat}\Bfootnote{%
\textit{(1)}~$L\lambda$ esse triangulis %
\textit{(2)}~rectas~\textit{L}}}
$L\lambda$ esse triangulis $CB\beta$ proportionales, seu 
\edtext{triangulis $CB\beta$ et $C(B)(\beta)$}{\lemma{triangulis}\Bfootnote{%
\textit{(1)}~existentibus aequalibus, e %
\textit{(2)}~$CB\beta$ et $C(B)(\beta)$~\textit{L}}}
existentibus aequalibus, etiam rectas $L\lambda$ et 
\edtext{$(L)\lbrack(\lambda)\rbrack$}{\lemma{}\Bfootnote{$(L)\lambda$ \textit{L ändert Hrsg.}}}
esse aequales, quod quidem in circulo verissimum est, necesse est centrum gravitatis\indexS{centrum gravitatis semicircumferentiae} 
\edtext{omnium $L\lambda$ seu semicircumferentiae}{\lemma{omnium $L\lambda$}\Bfootnote{%
\textit{(1)}~idem esse cu %
\textit{(2)}~seu %
\textit{(a)}~semicirculi %
\textit{(b)}~semicircumferentiae~\textit{L}}}
\textit{FeD} idem esse cum centro gravitatis semicirculi\indexS{centrum gravitatis semicirculi} \textit{ABD}. 
\edtext{%
Centrum enim 
\edtext{gravitatis \bracks{\textit{e}} cujuslibet $L\lambda$ est}{\lemma{gravitatis}\Bfootnote{\bars{\textit{E} \textit{erg., ändert Hrsg.}} %
\textit{(1)}~omnium $L\lambda$, sunt c %
\textit{(2)}~cujuslibet $L\lambda$ est~\textit{L}}}
centrum gravitatis cujuslibet $CB\beta$, et si cum $CB\beta$ sint ipsis $L\lambda$ proportionales nihil refert ad centrum gravitatis\indexS{centrum gravitatis}
\edtext{commune}{\lemma{}\Bfootnote{commune \textit{erg.~L}}} 
omnium ponderum\indexS{centrum gravitatis commune omnium ponderum} ex suis centris
gravitatis particularibus 
\edtext{\bracks{\textit{e}}}{\lemma{}\Bfootnote{\textit{E} \textit{L ändert Hrsg.}}} 
suspensorum utrum 
\edtext{pondera illa sint}{\lemma{pondera}\Bfootnote{%
\textit{(1)}~inde s %
\textit{(2)}~illa sint~\textit{L}}}
recta, an 
\edtext{ipsis proportionalia triangula. Idem}{\lemma{ipsis}\Bfootnote{%
\textit{(1)}~proportionales %
\textit{(2)}~proportionalia triangula.
\textit{(a)}~Centrum igitur gravitatis %
\textit{(b)}~Idem~\textit{L}}}
argumentum procedit etiam in sectore quocunque. Itaque 
\textls{centrum gravitatis semicirculi\indexS{centrum gravitatis semicirculi} vel sectoris cujuscunque\indexS{centrum gravitatis sectoris circuli cujuscunque} 
idem est cum centro gravitatis semicircumferentiae\indexS{centrum gravitatis semicircumferentiae}}%
}{%
\applabel{FnA:37_05_064-065_elementa-lineae}%
\lemma{}%
\Afootnote{\textit{Am Rand}: 
Si \indexS{elementum loci centrorum gravitatis}elementa\textsuperscript{\bracks{a}} 
lineae quae est locus centrorum gravitatis\indexS{locus centrorum gravitatis}\indexS{centrum gravitatis}
sint proportionalia ipsis quorum sunt illa centra, 
centrum gravitatis loci\indexS{centrum gravitatis loci centrorum gravitatis} 
erit centrum gravitatis totale\indexS{centrum gravitatis totale}
eorum quorum centra gravitatis in illum locum cadunt. 
\\[1.5ex] %
{\footnotesize 
\textsuperscript{\bracks{a}} elementa %
\textit{(1)}~loci g %
\textit{(2)}~lineae %
\textit{(a)}~quae %
\textit{(b)}~quae est locus~\textit{L}}}} 
\edtext{\textls{vel arcus\indexS{centrum gravitatis arcus} alterius}}{\lemma{\textls{vel}}\Bfootnote{%
\textit{(1)}~\textls{similis} %
\textit{(2)}~\textls{arcus} %
\textit{(a)}~\textls{similis sectoris} %
\textit{(b)}~sectoris in %
\textit{(aa)}~sesquia %
\textit{(bb)}~subsesquialtera 
\textit{(b)}~\bars{similis \textit{erg.}} et \textls{similiter positi (seu concentrici)} sectoris in subsesquialtera ratione minoris %
\textit{(3)}~\textls{alterius}~\textit{L}}}
\textls{sectoris\indexS{centrum gravitatis sectoris concentrici} similis et similiter positi seu}
\edtext{\textls{concentrici, per radium}}{%
\lemma{\textls{concentrici,}}%
\Bfootnote{%
\textit{(1)}~radio
\textit{(2)}~a radio
\textit{(3)}~\textls{per radium}~\textit{L}}}
\textls{in subsesquialtera ratione minorem descripti}.
\pend
\pstart 
Simili plane ratione 
\edtext{\textls{Hemisphaerii vel}}{\lemma{\textls{Hemisphaerii}}\Bfootnote{%
\textit{(1)}~\textls{et} %
\textit{(2)}~\textls{vel}~\textit{L}}}
\textls{sectoris solidi, centrum gravitatis\indexS{centrum gravitatis Hemisphaerii}\indexS{centrum gravitatis sectoris solidi} coincidit cum centro gravitatis superficiei hemisphaerii\indexS{centrum gravitatis superficiei hemisphaerii} vel sectoris minoris similis et similiter positi in subsesquitertia ratione minoris}, 
quoniam Haemisphaerium vel ejus sector in innumera conformia corpora resolvi\indexS{resolutio Haemisphaerii in innumera corpora conformia} 
potest, quorum apex in centro basis in superficie, centrum gravitatis\indexS{centrum gravitatis} autem abest a 
\edtext{centro tribus quartis radii. Sunt ergo}{\lemma{centro}\Bfootnote{%
\textit{(1)}~$\frac{3}{4}$ rad %
\textit{(2)}~tribus
\textit{(a)}~quatris
\textit{(b)}~quartis radii. %
\textit{(aa)}~Unde %
\textit{(bb)}~Sunt %
\textit{(aaa)}~autem %
\textit{(bbb)}~ergo~\textit{L}}}
loca omnium centrorum\indexS{locum omnium centrorum gravitatis} partium sectoris in minoris sectoris superficie; et 
\edtext{spatia seu partes superficiei sectoris}{\lemma{spatia}\Bfootnote{%
\textit{(1)}~inter %
\textit{(2)}~in
\textit{(3)}~seu partes superficiei %
\textit{(a)}~hemishpaerii %
\textit{(b)}~he
\textit{(c)}~sectoris~\textit{L}}}
minoris, a majore sectore abscissae, sunt proportionales ipsis majoribus sectoribus, et eadem habere intelligi possunt centra gravitatis\indexS{centrum gravitatis} cum ipsis sectoribus. 
\edtext{Ergo nihil}{%
\lemma{Ergo}%
\Bfootnote{%
\textit{(1)}~cum
\textit{(2)}~nihil~\textit{L}}}
refert ad centrum gravitatis totale,\indexS{centrum gravitatis totale} 
utrum ex his centris partialibus suspendas 
\edtext{partes sectoris}{\lemma{partes}\Bfootnote{%
\textit{(1)}~hemisphaerii %
\textit{(2)}~sectoris~\textit{L}}}
an partes respondentes superficiei sectoris minoris. Ac proinde centrum gravitatis\indexS{centrum gravitatis}
erit idem utrobique. 
\edtext{Si in Hemisphaerii}{\lemma{Si}\Bfootnote{%
\textit{(1)}~Hemisphaeri %
\textit{(2)}~in Hemisphaerii~\textit{L}}}
majoris 
\edtext{superficie circulos}{\lemma{superficie}\Bfootnote{%
\textit{(1)}~majores %
\textit{(2)}~circulos~\textit{L}}}
inassignabiles descriptos 
\edtext{intelligas, sese tangentes,}{\lemma{intelligas,}\Bfootnote{%
\textit{(1)}~qui sese %
\textit{(2)}~quorum tres quilibet sese tan %
\textit{(3)}~sese tangentes,~\textit{L}}}
et tam conos inde ortos quam intercepta inter hos conos corpora conicos consideres 
\edtext{circulosque similiter}{\lemma{circulosque}\Bfootnote{%
\textit{(1)}~similes %
\textit{(2)}~similiter~\textit{L}}}
positos et proportionales, interceptaque his 
\edtext{spatia trilinea (si vis tresquosque circulos se tangere)}{\lemma{spatia}\Bfootnote{%
\textit{(1)}~in su %
\textit{(2)}~in quas superficies minoris %
\textit{(3)}~triang %
\textit{(4)}~tri %
\textit{(5)}~trilinea %
\textit{(a)}~si placet aut %
\textit{(aa)}~ali %
\textit{(bb)}~quadrilinea (prout circul %
\textit{(b)}~(si vis %
\textit{(aa)}~circulos in %
\textit{(bb)}~tresquosque circulos \bars{aequales \textit{gestr.}} se tangere)~\textit{L}}}
vel quadrilinea aliave, consideres, res est manifesta.
\pend
\clearpage
\pstart\noindent
\href{\Scan{LH_37_05_064v}}%
{\lbrack64~v\textsuperscript{o}\rbrack}  
\pend
\nobreak	%
\begin{edfigure}[0.5em][1em]
\includegraphics[scale=1]{%
gesamttex/edit_VIII-4/images/LH_37_05_064-065_d2_064v.pdf}
\edcaption{\edtext{\lbrack\textit{Fig.~2}\rbrack}{%
\lemma{\lbrack\textit{Fig.~2}\rbrack}%
\Cfootnote[nonum]{Zwei gestrichene Entwürfe zum Diagramm werden nicht wiedergegeben.}}}
\end{edfigure}
\pstart
Itaque haec quidem in circulo pulchre procedunt sed difficulter ad alias figuras
\edtext{transferuntur, patet tamen si figura aliqua plana}{\lemma{transferuntur,}\Bfootnote{%
\textit{(1)}~si scilicet %
\textit{(2)}~patet tamen si figura aliqua %
\textit{(a)}~ita resolvi poss %
\textit{(b)}~plana~\textit{L}}}
ita resolvi\indexS{resolutio figurae in partes inassignabiles} posset in
\edtext{partes inassignabiles\indexS{partes inassignabiles figurae} ut quadrilatera \textit{\ind{1}L\ind{1}M\ind{2}M\ind{2}L} ut linea}{\lemma{partes}\Bfootnote{%
\textit{(1)}~\textit{ML}, verbi grat %
\textit{(2)}~ut %
\textit{(3)}~inassignabiles ut quadrilatera \textit{\ind{1}L\ind{1}M\ind{2}M\ind{2}L} %
\textit{(a)}~\bars{ut \textit{streicht Hrsg.}} centra gravitatis %
\textit{(b)}~ut %
\textit{(aa)}~du %
\textit{(bb)}~loco %
\textit{(cc)}~linea~\textit{L}}}
\textit{EE} ducta, quae sit locus centrorum gravitatis\indexS{locus centrorum gravitatis} \textit{E}, 
harum partium sint elementa hujus curvae ut \textit{\ind{1}E\ind{2}E} ipsis partibus in quas figura plana resolvitur proportionalia; 
eodem plane modo manifestum erit, centrum gravitatis curvae\indexS{centrum gravitatis curvae} \textit{EE}, et figurae \textit{LLLMMM} coincidere. 
Praeterea patet etiam ex dimensione curvae \textit{EE} dimensionem figurae illius haberi posse. 
Tale quid etiam in figuris solidis tentari potest, ubi major adhuc est eligendi libertas. 
\pend
\pstart
Videamus an haec cum fructu aliquo transferri possint
ad curvas evolutione descriptas,\indexS{curva evolutione descripta}
\edtext{quanquam jam tum circa illas talia quaedam olim me inquisivisse meminerim.}{%
\lemma{quanquam \bracks{...}  meminerim.}%
\Cfootnote{Womöglich sind hiermit die vermutlich Aufzeichnungen N.\;\ref{Cc2-1421} gemeint, insbesondere der zweite Text: \refpassage{35_08_30_133-teil2a}{35_08_30_133-teil2b}}}
\pend
\begin{edfigure}%
\includegraphics[scale=1]{%
gesamttex/edit_VIII-4/images/LH_37_05_064-065_d3_064v.pdf}
\edcaption[0.25em]{\lbrack\textit{Fig.~3}\rbrack}
\end{edfigure}
\pstart 
Evolutione curvae vel polygoni 
\edtext{inassignabilis}{\lemma{}\Bfootnote{inassignabilis \textit{erg.~L}}}
\textit{AMB} describatur curva vel polygonum inassignabile \textit{ALC}, et spatium \textit{AMBCLA} resolutum erit in triangula \textit{A\ind{1}M\ind{1}L}, \textit{\ind{1}L\ind{2}M\ind{2}L}, etc. quorum centra gravitatis\indexS{centra gravitatis triangulorum} sint \textit{\ind{1}E}, \textit{\ind{2}E}, \textit{\ind{3}E}, etc. Si jam 
\edtext{contingeret elementa}{\lemma{contingeret}\Bfootnote{%
\textit{(1)}~intervalla %
\textit{(2)}~elementa~\textit{L}}}
curvae \textit{EE}, seu intervalla duorum centrorum gravitatis\indexS{intervallum duorum centrorum gravitatis} esse triangulis porportionalia (vel quod eodem
\edtext{rediret inassignabilibus)}{%
\lemma{rediret}%
\Bfootnote{%
\textit{(1)}~\bars{in \textit{streicht Hrsg.}}
\textit{(2)}~in
\textit{(3)}~inassignabilibus)~\textit{L}}}
duorum in quae cadunt triangulorum summae (tantum enim omnia Triangula praeter duo extrema quae nihil addunt vel tollunt, duplicata fingenda essent) foret idem centrum gravitatis curvae\indexS{centrum gravitatis curvae evolutione descriptae} \textit{EE}, et figurae 
\edtext{ex}{\lemma{}\Bfootnote{ex \textit{erg.~L}}} 
his triangulis compositae.\indexS{centrum gravitatis figurae triangulis compositae} 
Equidem triangula sunt in duplicata ratione rectarum \textit{ML}, quae vocabo
\textit{m} et \textit{ME} 
\edtext{est \rule[-3.5mm]{0pt}{9mm}$\displaystyle\frac{2}{3}m$ et}{\lemma{est $\displaystyle\frac{2}{3}m$}\Bfootnote{%
\textit{(1)}~et \textit{\ind{1}E\ind{2}E} est 
\textit{(a)}~\textit{d}
\textit{(b)}~\textit{d}
\textit{(c)}~$\overline{\displaystyle\frac{2}{3}d\omega}^2$ %
\textit{(2)}~et~\textit{L}}}
\textit{\ind{1}L\ind{2}L} sit $\displaystyle\frac{m}{a}\omega$. Posito $\omega$ esse inassignabilem 
\edtext{constantem, erit \textit{\ind{1}E\ind{2}E}}{\lemma{constantem,}\Bfootnote{%
\textit{(1)}~erit %
\textit{(2)}~erit %
\textit{(a)}~\textit{\ind{1}E\ind{2}E} aeq. \textit{d} %
\textit{(b)}~\textit{\ind{1}E\ind{2}E}~\textit{L}}}
aeq. $\sqrt{\overline{\displaystyle\frac{2}{3}dm}^2+\overline{\displaystyle\frac{2}{3}\displaystyle\frac{m}{a}d\omega}^2}$, 
\edtext{quae si}{\lemma{quae}\Bfootnote{%
\textit{(1)}~debet esse ad $\omega$ %
\textit{(2)}~si~\textit{L}}}
debet esse triangulo seu ipsi \textit{mm} proportionalis fiet: 
aeq. $d\omega \displaystyle\frac{mm}{aa}$. 
Ergo fiet: \rule[-5mm]{0pt}{11mm}$\displaystyle\frac{4}{9}d\overline{m}^2+\displaystyle\frac{4mm}{9aa}d\overline{\omega}^2$ aeq. \rule[-5mm]{0pt}{11mm}$\displaystyle\frac{m^4}{a^4}d\overline{\omega}^2$, seu fiet: $d\omega$ aeq. \rule[-5mm]{0pt}{11mm}$\displaystyle\frac{2aadm}{m\sqrt{9m^2-4aa}}$. 
\pend
\pstart
Itaque si curva data \textit{AMMB} ita secetur in partes, 
\edtext{ut elementis ejus existentibus \textit{dm} fiat}{\lemma{ut}\Bfootnote{%
\textit{(1)}~elementa ejus respectu alicujus inassignabilis $d\omega$, tum ordinatim assumtae %
\textit{(2)}~elementis ejus existentibus 
\textit{(a)}~\textit{dm},
\textit{(b)}~\textit{dm}
\textit{(aa)}~eductisque inde rectis utcunque, ali %
\textit{(bb)}~fiat~\textit{L}}}
$\displaystyle\frac{2aadm}{m\sqrt{9mm-4aa}}$\rule[-5mm]{0pt}{11mm} quantitas constans\bracks{,} tunc proportionales erunt ipsae \textit{\ind{1}E\ind{2}E}, ipsis triangulis \textit{\ind{1}L\ind{2}M\ind{2}L}\bracks{,} 
quantitas autem talis constans pendet a quadratura seu summa ipsius \rule[-5mm]{0pt}{11mm}$\displaystyle\frac{2aadm}{m\sqrt{9mm-4aa}}$, et 
\edtext{talis sectio a curva}{\lemma{}\Bfootnote{%
talis \bars{figurae \textit{gestr.}} sectio
\textit{(1)}~pendet %
\textit{(2)}~a curva~\textit{L}}}
\textit{AMM} pendet a dimensione ipsius et partium ejus\bracks{,} ita enim abscindi poterit arcus in data 
\edtext{relatione. Sed}{\lemma{relatione.}\Bfootnote{%
\textit{(1)}~\bars{Sed \textit{streicht Hrsg.}} videtur tali sectio %
\textit{(2)}~Sed~\textit{L}}}
ecce, 
\edlabel{37_05_064-065_4a}%
nonne videtur tali sectione actuali non esse opus, sufficereque 
\edtext{ipsam esse}{\lemma{}\Bfootnote{ipsam \bars{ipsam \textit{streicht Hrsg.}} esse~\textit{L}}}
possibilem\bracks{,}%
\edlabel{37_05_064-065_4b}%
si enim secundum illam sectionem centrum gravitatis curvae\indexS{centrum gravitatis curvae} \textit{EE}, et figurae\indexS{centrum gravitatis figurae} ubique coincidit, secundum quamcumque coincidet. Itaque hic videtur subesse subtile sophisma. 
\pend 
\pstart\noindent
\href{\Scan{LH_37_05_065r}}%
{\lbrack65~r\textsuperscript{o}\rbrack} 
\pend
\nobreak	%
\begin{edfigure}%
\includegraphics[scale=1]{%
gesamttex/edit_VIII-4/images/LH_37_05_064-065_d4_065r.pdf}
\edcaption{\lbrack\textit{Fig.~4}\rbrack}
\end{edfigure}
\pstart
Similis difficultas tentetur in simpliciori casu. Propos\angles{ita cu}rva \textit{ALC}, 
radiis 
\edtext{ex puncto}{\lemma{ex}\Bfootnote{%
\textit{(1)}~centro %
\textit{(2)}~puncto~\textit{L}}}
\textit{B} eductis secetur arcus in triangula ut \textit{B\ind{1}L\ind{2}L}, sic ut sumtis \textit{BE} in d\angles{a}ta quadam ratione constante, (vel si mavis relatione) ad \textit{BL}, 
\edtext{sit \textit{\ind{1}E\ind{2}E},}{\lemma{sit}\Bfootnote{%
\textit{(1)}~ipsa %
\textit{(2)}~el %
\textit{(3)}~\textit{\ind{1}E\ind{2}E},~\textit{L}}}
elementum curvae\indexS{elementum curvae} \textit{EE}, ipsi triangulo \textit{B\ind{1}L\ind{2}L} proportionale. Quod quidem si possibile esset, ostendi posset curvam illam \textit{EE}, fore spatiis respondentibus proportionalem, quod tamen cum data quaevis curva pro \textit{EE}, assumi queat, quia generaliter de quacunque relatione ipsius \textit{BE} ad \textit{BL} ratiocinationem instituimus, fieri non potest. Rationem 
\edtext{tamen impossibilitatis}{\lemma{tamen}\Bfootnote{%
\textit{(1)}~absurditatis %
\textit{(2)}~impossibilitatis~\textit{L}}}
a priori nondum video, sed deprehenderem haud dubie si methodum exequi aggrederer.
\pend
\pstart
Aliud autem non 
\edtext{inelegans hic}{\lemma{inelegans}\Bfootnote{%
\textit{(1)}~. Sumatur %
\textit{(2)}~hic~\textit{L}}}
incidit. Sumatur in \textit{B\ind{1}L} ipsa \textit{B\ind{1}V} aequalis ipsi \textit{B\ind{2}L}, et ita similiter in caeteris, ut designentur puncta \textit{\ind{1}V}, \textit{\ind{2}V}, \textit{\ind{3}V}, manifestum est triangulum \textit{B\ind{1}L\ind{2}L} seu \textit{B\ind{1}V\ind{2}L},
\edtext{(nam Anoristis\indexS{Anoristum sive inassignabile} sive inassignabilibus discrimen anoristum\indexS{discrimen anoristum}}{\lemma{(nam}\Bfootnote{%
\textit{(1)}~assign %
\textit{(2)}~Anoristicum %
\textit{(3)}~Anoristis sive inassignabilibus discrimen
\textit{(a)}~A
\textit{(b)}~anoristum~\textit{L}}}
est\lbrack)\rbrack. Porro cum ipsae \textit{BL}, sint perpendiculares vel quasi ad ipsas \textit{\ind{1}V\ind{2}L}, hinc si ex \textit{L} in axem et basin demittantur 
\edtext{perpendiculares \textit{LN},}{\lemma{perpendiculares}\Bfootnote{%
\textit{(1)}~ut \textit{\ind{1}L\ind{2}N}, $\ind{2}$ %
\textit{(2)}~\textit{LN},~\textit{L}}}
et \textit{LQ}, patet omnia, \textit{NL} in \textit{\ind{1}V\ind{2}L}, 
\edtext{esse momentum}{\lemma{esse}\Bfootnote{%
\textit{(1)}~momenta omn %
\textit{(2)}~momentum~\textit{L}}}
omnium \textit{\ind{1}V\ind{2}L} ex axe, et omnia, \textit{LQ} in \textit{\ind{1}V\ind{2}L} esse momentum omnium \textit{\ind{1}V\ind{2}L} ex basi. Jam \textit{BL} in \textit{\ind{1}N\ind{2}N} aeq. \textit{NL} in \textit{\ind{1}V\ind{2}L} et similiter \textit{BL} in \textit{\ind{1}Q\ind{2}Q} aeq. \textit{QL} in \textit{\ind{1}V\ind{2}L}. Ergo radii ad 
\edtext{axem applicati dant momentum}{\lemma{axem}\Bfootnote{%
\textit{(1)}~sunt %
\textit{(2)}~applicati %
\textit{(a)}~sunt %
\textit{(b)}~dant %
\textit{(aa)}~momenta %
\textit{(bb)}~momentum~\textit{L}}}
differentiarum radiorum, ex axe, et radii ad basin applicati dant momentum\indexS{momentum differentiarum radiorum} differentiarum radiorum ex basi, et si ratio inter duas illas figuras detur, dabitur et recta transiens per centrum gravitatis harum differentiarum. \textit{BL} vocemus \textit{l} et \textit{\ind{1}V\ind{2}L} vocemus \textit{dv}. Fiet \textit{dvl} triangulum \textit{B\ind{1}L\ind{2}L}, et $\int ldv$ aeq. areae figurae. 
Porro quaeratur $\theta$ aeq. 
\rule[-3mm]{0pt}{8mm}%
\edtext{$\displaystyle\frac{ldv}{a}$. Erit}{\lemma{$\displaystyle\frac{ldv}{a}$}\Bfootnote{%
\textit{(1)}~quae %
\textit{(2)}~.~Erit~\textit{L}}}
$\theta$ proportionalis triangulo, ea $\theta$ sit in figura $\theta\mu$, ut $\ind{1}\theta\ind{1}\mu$, $\ind{2}\theta\ind{2}\mu$, parallela ipsi \textit{\ind{1}V\ind{2}L}, \textit{\ind{2}V\ind{3}L} collocata, patet ipsas $\ind{1}\theta\ind{1}\mu$, $\ind{2}\theta\ind{2}\mu$ fore triangulis \textit{B\ind{1}L\ind{2}L}, \textit{B\ind{2}L\ind{3}L}, etc. proportionales, jam 
\edtext{$B\ind{1}\theta$ est \textit{B\ind{1}L}, seu ad \textit{B\ind{1}V} seu}{\lemma{$B\ind{1}\theta$}\Bfootnote{%
\textit{(1)}~sit %
\textit{(2)}~est ad \textit{B\ind{1}L},
\textit{(a)}~ut %
\textit{(b)}~ut %
\textit{(c)}~seu
ad \textit{B\ind{1}V} 
\textit{(aa)}~ut
\textit{(bb)}~seu~\textit{L}}}
ad \textit{l}, ut $\theta$ ad \textit{dv} seu ut \textit{l} ad \textit{a}, ergo $B\ind{1}\theta$ aeq. \rule[-3mm]{0pt}{8mm}$\displaystyle\frac{ll}{a}$, hinc ipsa $B\ind{1}\theta$ \rule[-3mm]{0pt}{8mm}%
\edtext{seu $\displaystyle\frac{ll}{a}$ applicata}{\lemma{seu}\Bfootnote{%
\textit{(1)}~\textit{lla} %
\textit{(2)}~$\displaystyle\frac{ll}{a}$ %
\textit{(a)}~transferens %
\textit{(b)}~applicata~\textit{L}}}
ad axem erunt momentum\indexS{momentum} ipsorum $\theta\mu$ ex axe, translata
ad basin erunt momentum\indexS{momentum} ipsorum $\theta\mu$ ex basi, et datis quadraturis translatorum tam ad 
\edtext{ax\lbrack em\rbrack}{\lemma{}\Bfootnote{axin \textit{L ändert Hrsg.}}}
quam ad basin, vel saltem eorum ratione, datur recta ex centro \textit{B} educta transiens per centrum gravitatis\indexS{centrum gravitatis} ipsorum $\theta\mu$ omnium. %
\pend
\pstart
Porro si ponamus in \textit{\ind{1}E} esse centrum gravitatis trianguli\indexS{centrum gravitatis trianguli} \textit{B\ind{1}L\ind{2}L}, erit \textit{B\ind{1}E} aequ. $\displaystyle\frac{2}{3}l$. Ergo $B\ind{1}\theta$ ad \textit{B\ind{1}E} ut $\displaystyle\frac{ll}{a}$ ad \rule[-4mm]{0pt}{8mm}$\displaystyle\frac{2}{3}l$ seu ut \textit{l} ad $\displaystyle\frac{2}{3}a$. 
\pend
\pstart
Video tamen etsi recta habeatur quae momentum ipsarum $\theta\mu$ proportionalium  ipsius \textit{B\ind{1}L\ind{2}L} bisecat, non tamen ipsorum triangulorum \textit{B\ind{1}L\ind{2}L} momentum ab ea recta bisecari, quia 
\edtext{non sunt sita}{%
\lemma{non sunt}%
\Bfootnote{%
\textit{(1)}~proportionaliter
\textit{(2)}~sita~\textit{L}}}
triangula similiter ut ipsa $\theta\mu$, nam $B\theta$ sunt ut \textit{ll}, 
\edtext{et \textit{BL} seu \textit{BE}}{%
\lemma{et \textit{BL}}%
\Bfootnote{%
\textit{(1)}~sunt
\textit{(2)}~seu \textit{BE}~\textit{L}}}
sunt ut \textit{l}, 
\edtext{hinc rectae jungentis}{%
\lemma{hinc}%
\Bfootnote{%
\textit{(1)}~recta jungens
\textit{(2)}~rectae jungentis~\textit{L}}}
duorum $\theta\mu$ centra medium non cadit necessario 
\edtext{in rectae jungentis duo}{\lemma{in rectae jungentis}\Bfootnote{%
\textit{(1)}~duorum %
\textit{(2)}~duo~\textit{L}}}
\textit{E} (quae sunt triangulorum centra) medium. 
\pend
\begin{edfigure}%
\includegraphics[scale=1]{%
gesamttex/edit_VIII-4/images/LH_37_05_064-065_d5_065r.pdf}
\edcaption{\lbrack\textit{Fig.~5}\rbrack}
\end{edfigure}
\pstart
Illud superest
\edtext{considerandum}{\lemma{}\Bfootnote{considerandum \textit{erg.~L}}}
an momentum\indexS{momentum figurae} alicujus figurae sumere liceat ex axe curvo,\indexS{axis curvus} 
\edtext{v.\,g. an}{\lemma{v.\,g.}\Bfootnote{%
\textit{(1)}~momentum %
\textit{(2)}~an~\textit{L}}}
momentum figurae\indexS{momentum figurae} \textit{ABC} 
\edtext{ex linea}{\lemma{ex}\Bfootnote{%
\textit{(1)}~axe %
\textit{(2)}~linea~\textit{L}}}
curva \textit{ALLC} intelligi possit, hoc est 
\edtext{eductis curvae perpendicularibus \textit{LM}, usque ad axem \textit{AM}, intelligatur quadrilaterum}{\lemma{eductis}\Bfootnote{%
\textit{(1)}~ex curva perpendicularibus \textit{LM} ad axem %
\textit{(2)}~curvae perpendicularibus %
\textit{(a)}~usque a %
\textit{(b)}~\textit{LM} usque ad axem \textit{AM}, %
\textit{(aa)}~utrum %
\textit{(bb)}~intelligatur %
\textit{(aaa)}~triangulum %
\textit{(bbb)}~rectangulum %
\textit{(ccc)}~quadrilaterum~\textit{L}}}
\textit{\ind{1}M\ind{1}L\ind{2}L\ind{2}M} mobile circa \textit{\ind{1}L\ind{2}L} velut axem, 
\edtext{et quadrilaterum}{\lemma{et}\Bfootnote{%
\textit{(1)}~triangulum \bars{rectangulum \textit{erg.}} \textit{\ind{2}M\ind{2}L\ind{3}L} mobile circa %
\textit{(2)}~quadrilaterum~\textit{L}}}
\textit{\ind{2}M\ind{2}L\ind{3}L\ind{3}M} mobile circa \textit{\ind{2}L\ind{3}L} velut axem.
Patet his rotationibus (quanquam unaquaeque separatim fieri concipienda sit)
fieri solidum aliquod cujus area habebitur 
\edtext{si centri gravitatis ipsius quadrilateri distantia ab axe}{\lemma{si}\Bfootnote{%
\textit{(1)}~centrum %
\textit{(2)}~centrorum %
\textit{(3)}~centri gravitatis %
\textit{(a)}~dista %
\textit{(b)}~ipsius quadrilateri distantia ab
\textit{(aa)}~element %
\textit{(bb)}~ab %
\textit{(cc)}~axe~\textit{L}}}
suo, hoc est a tangente curvae ducta in ipsum quadrilaterum\bracks{,}
ipsi curvae superponatur, talis 
\edtext{superficies ex}{\lemma{superficies}\Bfootnote{%
\textit{(1)}~cylindrica erit solido proportionalis %
\textit{(2)}~ex~\textit{L}}}
cylindrica quodammodo absecta, erit solido illi proportionalis. Quod 
\edtext{si in rectum}{\lemma{si}\Bfootnote{%
\textit{(1)}~solidum %
\textit{(2)}~in rectum~\textit{L}}}
extendamus curvam, et 
\edtext{tum applicare ei possumus}{\lemma{tum}\Bfootnote{%
\textit{(1)}~applicamus ei jam %
\textit{(2)}~applicare ei possumus~\textit{L}}}
rectangula, ipsis quadrilateris 
\edtext{aequalia vel}{\lemma{aequalia}\Bfootnote{%
\textit{(1)}~habentiaque %
\textit{(2)}~vel~\textit{L}}}
proportionalia, habentiaque etiam centra gravitatis sua aeque remota ab illa recta tanquam 
\edtext{axe, vel}{\lemma{axe,}\Bfootnote{%
\textit{(1)}~ut %
\textit{(2)}~vel~\textit{L}}}
proportionaliter, ut ante a tangente curvae, tunc rotatione inde factae curvae natum solidum\indexS{solidum rotatione curvae natum} foret proportionale priori. Sed tale quid communiter fieri non potest. Dubito an in 
\edtext{rotatione punctum \textit{\ind{2}M} non}{\lemma{rotatione}\Bfootnote{%
\textit{(1)}~puncta \textit{\ind{1}M} et \textit{\ind{2}M} %
\textit{(2)}~punctum \textit{\ind{2}M}
\textit{(a)}~prop
\textit{(b)}~non~\textit{L}}}
divelli debeat in duo. Imo utrum non solidum ab 
\edtext{unius quadrilateri}{\lemma{unius}\Bfootnote{%
\textit{(1)}~rectanguli %
\textit{(2)}~quadrilateri~\textit{L}}}
rotatione genitum\indexS{solidum ab rotatione quadrilateri genitum} 
vel in alterius proximi rotationem irrumpat, vel locum vacuum relinquat,
et ita solum hac ratione non completur nec proinde tale quid  
\edlabel{37_05_064-065_2a}%
\edtext{}{\xxref{37_05_064-065_2a}{37_05_064-065_2b}\lemma{procedit.}\Bfootnote{%
\href{\Scan{LH_37_05_065v}}%
{\lbrack65~v\textsuperscript{o}\rbrack} 
\textit{(1)}~Curvam ita secare, \bars{ut el \textit{streicht Hrsg.}}
\textit{(2)}~Si velimus~\textit{L}}}%
procedit.
\pend
\clearpage
\pstart
\noindent
\href{\Scan{LH_37_05_065v}}%
{\lbrack65~v\textsuperscript{o}\rbrack}
\pend
\nobreak
\begin{edfigure}[0.5em][1em]
\includegraphics[scale=1]{%
gesamttex/edit_VIII-4/images/LH_37_05_064-065_d6_065v.pdf}
\edcaption{\edtext{\lbrack\textit{Fig.~6}\rbrack}{%
\lemma{\lbrack\textit{Fig.~6}\rbrack}%
\Cfootnote[nonum]{Ein gestrichener Entwurf zum Diagramm wird nicht wiedergegeben.}}}
\end{edfigure}
\pstart
Si velimus\edlabel{37_05_064-065_2b} 
ut elementa lineae centrorum gravitatis\indexS{elementa lineae centrorum gravitatis rectangulorum}
omnium rectangulorum componentium sint rectangulis 
\edtext{proportionali\lbrack a\rbrack}{\lemma{}\Bfootnote{proportionalis \textit{L ändert Hrsg.}}} 
idem est ac 
\edtext{si jungamus}{\lemma{si}\Bfootnote{%
\textit{(1)}~faciamus %
\textit{(2)}~jungamus~\textit{L}}}
dimidia puncta ordinatarum et ita faciamus novam curvam, 
\edtext{quae ita}{\lemma{}\Bfootnote{quae \bars{quod \textit{streicht Hrsg.}} ita~\textit{L}}}
secanda. Sed cum elementa hujus novae curvae sint proportionalia elementis curvae\indexS{elementum curvae} datae, 
\edtext{%
(nam datae Elementa sunt 
\edtext{$\sqrt{d\overline{x}^2+d\overline{y}^2}$\lbrack,\rbrack\ novae}{\lemma{$\sqrt{d\overline{x}^2+d\overline{y}^2}$}\Bfootnote{%
\textit{(1)}~quae sunt %
\textit{(2)}~novae~\textit{L}}}
sunt 
\rule[-4mm]{0pt}{10mm}%
$\sqrt{d\overline{\overline{\displaystyle\frac{1}{2}x}^2}+d\overline{\overline{\displaystyle\frac{1}{2}y}^2}}$, 
seu \rule[-4mm]{0pt}{10mm}%
$\displaystyle\frac{1}{2}\sqrt{d\overline{x}^2+d\overline{y}^2}$ \lbrack)\rbrack%
}{\lemma{}\Afootnote{\textit{Nebenrechnung am oberen Blattrand}: $\sqrt{d\overline{x}^2+d\overline{y}^2}$.\enskip $\surd \displaystyle\frac{1}{2} d\overline{x}^2+$ \lbrack\textit{Rechnung bricht ab.}\rbrack}} 
patet esse proportionalia. Sufficit itaque quaeri modum ita secandi curvam per parallelas ut elementa curvae et elementa figurae sint proportionalia.
\pend
\pstart 
Seu erit \textit{ydx} aequ. $a\sqrt{d\overline{x}^2+d\overline{y}^2}$, 
seu $y^2d\overline{x}^2$ aeq. $aadx^2+a^2dy^2$ seu 
\rule[-3mm]{0pt}{3mm}%
$\displaystyle\frac{y^2-a^2}{a^2}$ aeq. 
\rule[-3mm]{0pt}{3mm}%
$\displaystyle\frac{dy^2}{dx^2}$. seu 
\rule[-3mm]{0pt}{6mm}%
$dx\ \overset{\astrosun}{\text{aeq.}}\ \dfrac{ady}{\sqrt{yy-aa}}$. 
Seu \textit{y} sic explicanda est, per 
\edtext{aliam assumtam,}{%
\lemma{aliam}%
\Bfootnote{%
\textit{(1)}~,~ut
\textit{(2)}~assumtam,~\textit{L}}}
ut si pro \textit{dx} substituas valorem quem postulat relatio inter \textit{x} et \textit{y}, debet idem prodire, quod ex 
\rule[-3mm]{0pt}{6mm}%
$\displaystyle\frac{ady}{\sqrt{yy-aa}}$. 
Quod si aequatio \astrosun\ posset revocari ad communem, res credo adhuc foret facilior. 
\edlabel{37_05_064-065_3a}%
\edtext{}{%
\xxref{37_05_064-065_3a}{37_05_064-065_3b}%
\lemma{}\Afootnote{\textit{Am Rand}: Latet in his aliquid non contemnendum.}}%
Puto pendere \astrosun\ ex quadratura Hyperbolae.\indexS{quadratura Hyperbolae}
\pend
\pstart
\edtext{Possunt}{\lemma{Hyperbolae.}\Bfootnote{%
\lbrack/\rbrack\ \textit{(1)}~Quid %
\textit{(2)}~Potest %
\textit{(3)}~Possunt~\textit{L}}}
enim innumeris 
\edtext{modis dividi}{\lemma{modis}\Bfootnote{%
\textit{(1)}~fieri %
\textit{(2)}~dividi~\textit{L}}}
figurae, et quaeri.\edlabel{37_05_064-065_3b}
\pend
\begin{edfigure}%
\includegraphics[scale=1]{%
gesamttex/edit_VIII-4/images/LH_37_05_064-065_d7_065v.pdf}
\edcaption{\lbrack\textit{Fig.~7}\rbrack}
\end{edfigure}
\pstart
Videndum tamen an non aequatio \astrosun\ sit impossibilis, nisi in eo casu quo curva jam est determinata. Patet etiam talem quadratricem Hyperbolicam cujus aequatio est 
\rule[-3mm]{0pt}{8mm}%
$dx\ \text{aeq.}\ \displaystyle\frac{ady}{\sqrt{yy-aa}}$ 
habere hoc privilegium ut curvae elementa\indexS{elementum curvae} sint proportionalia partibus figurae.
\pend 
